\documentclass[11pt]{amsart}

\usepackage[margin=1.1in]{geometry}
\usepackage{amsmath,amssymb,amsthm,mathtools}
\usepackage{enumitem}
\usepackage{hyperref}

\newtheorem{theorem}{Theorem}[section]
\newtheorem*{theorem*}{Theorem}
\newtheorem{proposition}[theorem]{Proposition}
\newtheorem{lemma}[theorem]{Lemma}
\newtheorem{corollary}[theorem]{Corollary}
\theoremstyle{definition}
\newtheorem{definition}[theorem]{Definition}
\theoremstyle{remark}
\newtheorem{remark}[theorem]{Remark}

\newcommand{\Z}{\mathbb{Z}}
\newcommand{\R}{\mathbb{R}}
\newcommand{\N}{\mathbb{N}}
\newcommand{\Zi}{\mathbb{Z}[i]}
\DeclareMathOperator{\Nm}{N}
\newcommand{\nm}[1]{\Nm(#1)}

\title[Lattice Points on Short Arcs of Circles]{Lattice Points on Short Arcs of Circles:\\ A Sublogarithmic Bound via Gaussian Prime Descent}
\author{Craig Chen}
\date{}

\begin{document}

\begin{abstract}
A classical theorem of Jarn\'{i}k bounds the number of integer lattice points on a convex arc of length $L$ by $O(L^{2/3})$, and this exponent cannot be improved for arcs of arbitrary length. We show that once an arc is constrained to the much shorter scale $L \le C\sqrt{R}$ on a circle of radius $R$ -- the scale at which curvature first competes with lattice spacing -- the trivial polynomial bound can be replaced by an essentially logarithmic one: the number of lattice points on such an arc is $O_C(\log R / \log\log R)$. The argument is arithmetic rather than convex-geometric. We encode lattice points on the circle $x^2+y^2=N$ as Gaussian integers of norm $N$, build a Vandermonde-type determinant on a window of consecutive points whose nonvanishing forces the points apart, and strengthen this separation using a prime-by-prime divisibility argument over a geometrically increasing stack of prime intervals. When that argument fails to find enough ``missing'' primes, it instead exhibits an actual prime divisor of $N$, along which the entire configuration of points \emph{descends} -- via the splitting behavior of primes in $\Z[i]$ -- to the same problem on a smaller circle, and the argument recurses. We describe this proof in detail; it has been formally verified, end to end, in Lean~4 with Mathlib.
\end{abstract}

\maketitle

\section{Introduction}

How many lattice points can lie on a short stretch of a circle? Gauss's classical question asks how many points of $\Z^2$ lie \emph{inside} the disc of radius $R$, and the error in the approximation $\#\{(x,y)\in\Z^2 : x^2+y^2\le R^2\}\approx \pi R^2$ remains, a century and a half later, only partially understood. The question we take up here is its older, more rigid cousin: not how many lattice points lie near a circle, but how many lie precisely \emph{on} one, restricted further to a short arc of it.

The basic obstruction is purely geometric. Three points on a circle of radius $R$, no two closer together than some distance $\delta$, span a triangle of area at least $\delta^2$ \emph{a priori}, but at most $O(\delta \cdot R)$ once one uses that the triangle's circumradius is $R$; balancing the two forces $\delta \gg R^{-1/3}$, so three lattice points on a circle cannot be arbitrarily close together. Jarn\'{i}k turned this kind of triangle inequality into a global theorem.

\begin{theorem*}[Jarn\'{i}k, 1926 \cite{Jarnik1926}]
A convex curve of length $L$ in $\R^2$ passes through at most $O(L^{2/3})$ integer lattice points.
\end{theorem*}

Applied to an arc of length $L$ on a circle of radius $R$, this gives a bound of shape $O(L^{2/3})$, and the exponent $2/3$ is sharp for arcs of length comparable to $R$ itself: there exist circles whose lattice points cluster densely enough that no better exponent is possible in general. What is \emph{not} pinned down by this convexity argument alone is what happens when the arc is much shorter than $R$ -- short enough that only the arithmetic of the particular radius $R$, rather than the curve's convexity, can be responsible for how many lattice points it carries.

This is the question we study. Write $N=R^2$, so that a lattice point on the circle of radius $R$ is a solution to $x^2+y^2=N$, equivalently a Gaussian integer $z=x+iy$ of norm $N$. We show that on any arc whose length is at most a fixed constant multiple of $\sqrt{R}$ -- already far shorter than the regime in which Jarn\'{i}k's exponent is sharp -- the trivial bound coming from convexity can be replaced by something close to a constant.

\begin{theorem*}[Main theorem, informal version]
Fix $C>0$. There is $R_0=R_0(C)$ so that for all $R\ge R_0$: any arc of length at most $C\sqrt{R}$ on the circle $x^2+y^2=R^2$ contains at most
\[
1168\cdot\frac{\log R}{\log\log R}
\]
lattice points.
\end{theorem*}

A precise statement, with the exact hypotheses needed to make ``arc'' and ``length'' formal, is Theorem~\ref{thm:main} below. The method has nothing to do with convexity; it is arithmetic from start to finish. The two facts about $\Z[i]$ that make it work are:

\begin{enumerate}[label=(\roman*)]
\item Distinct points on the circle $z\bar z = N$ generate a nonvanishing Vandermonde-type determinant, and the size of that determinant can be bounded below using only the multiplicativity of the norm on $\Z[i]$. This already forces some pair of points in any sufficiently long run to be far apart -- our first, crude separation bound.
\item Each rational prime dividing $N$ splits, is inert, or ramifies in $\Z[i]$ according to its residue mod $4$, and dividing every point in the configuration by such a prime factor produces a \emph{smaller} instance of exactly the same problem, on a circle of radius $\sqrt{N/p}$ or $\sqrt{N/p^2}$. The configuration \emph{descends}.
\end{enumerate}

The bulk of the work is in deciding, for a given radius, which of these two phenomena to invoke and how to bookkeep the cost of using them: (i) is strengthened by an auxiliary divisibility argument that asks whether several small primes \emph{fail} to divide $N$ (in which case (i) alone already gives a strong bound), and falls back on (ii) precisely when that fails -- which, by a Mertens-type counting argument over a geometric ladder of prime intervals, is exactly when (ii) is guaranteed to apply. Iterating (ii) finitely many times, with parameters tuned as functions of $R$, produces the stated $\log R/\log\log R$ bound.

\subsection*{A note on the formalization} Every theorem, lemma, and numerical constant in this report has been checked mechanically: the proof is formalized end to end in Lean~4, built on Mathlib, in the \texttt{GaussianChain} library accompanying this report. We have not attempted to optimize the constant $1168$, nor the auxiliary parameters that produce it; both are simply what falls out of an explicit, formally checked choice of parameters, and we expect them to be improvable by a more careful analysis. This report follows the structure of the formalization quite closely but is written for a mathematician, not a proof assistant, and we have suppressed routine bookkeeping (floor functions, the difference between an inequality and its rounded integer form, and similar) wherever it does not affect the mathematical content.

\subsection*{Organization} Section~\ref{sec:notation} fixes notation. Section~\ref{sec:main} states the main theorem precisely and comments on its hypotheses and constants. Section~\ref{sec:strategy} gives an overview of the proof, which we then carry out in four stages occupying Sections~\ref{sec:act1} through~\ref{sec:act4}: the basic determinant bound (Section~\ref{sec:act1}), the missing-prime strengthening and the Mertens estimates it depends on (Section~\ref{sec:act2}), prime descent (Section~\ref{sec:act3}), and the final assembly of the dichotomy into the asymptotic bound (Section~\ref{sec:act4}). We close with remarks on constants, the limits of the present method, and directions for future work in Section~\ref{sec:remarks}.

\section{Notation}\label{sec:notation}

We write $f \ll g$, or $f=O(g)$, to mean that there is a constant $c$ with $|f|\le c\,g$ for all values of the relevant parameters; a subscript, as in $f \ll_C g$, indicates that the implied constant is allowed to depend on $C$. Unless stated otherwise, all implied constants are absolute and effective, though we have made no effort to optimize them.

We identify $\Z^2$ with the ring of Gaussian integers $\Zi = \{a+bi : a,b\in\Z\}$ in the usual way, $z=x+iy$. We write $\bar z$ for complex conjugation and
\[
\nm{z} := z\bar z = x^2+y^2 \in \Z_{\ge 0}
\]
for the norm, which is completely multiplicative: $\nm{zw} = \nm{z}\nm{w}$. A lattice point on the circle $x^2+y^2=N$ is exactly a Gaussian integer of norm $N$. We write
\[
\nm{w-z}
\]
for the squared Euclidean distance between two such points, which we also call their \emph{squared chord length}.

A rational prime $p$ behaves in one of three ways in $\Z[i]$:
\begin{itemize}
\item if $p\equiv 1\pmod 4$, $p$ \emph{splits}: $p=\pi\bar\pi$ for a Gaussian prime $\pi$, by Fermat's two-square theorem;
\item if $p\equiv 3\pmod 4$, $p$ is \emph{inert}: $p$ remains prime in $\Z[i]$;
\item $p=2$ \emph{ramifies}: $2=-i(1+i)^2$.
\end{itemize}
We exclude $p=2$ from the descent argument throughout (Section~\ref{sec:act3}); it plays no special role in the final bound and could be folded into the missing-prime branch instead, but we have not pursued this.

For $m_0 < n$ and $N\in\N$ we write
\[
\Sigma(m_0,n) := \sum_{m_0 < p \le n} \frac{\log p}{p}, \qquad
\Sigma_{\mathrm{miss}}(N; m_0,n) := \sum_{\substack{m_0<p\le n\\ p\nmid N}} \frac{\log p}{p}
\]
for the usual weighted prime-counting sum over an interval and its restriction to primes \emph{not} dividing $N$ -- the ``missing'' primes relative to $N$. These are the central bookkeeping quantities of Section~\ref{sec:act2}.

\section{The main theorem}\label{sec:main}

To state the theorem we need to fix what we mean by an arc, in a form general enough to be proved by induction on the radius (and hence amenable to formalization), without committing in advance to a particular parametrization of the circle by angle or arclength.

\begin{definition}
Let $N\in\N$ and let $z_0,z_1,\dots,z_{M-1}\in\Z[i]$ be Gaussian integers of norm $N$. We say the sequence is \emph{ordered along an arc of width $L$} if there is a nondecreasing sequence $t_0\le t_1\le \cdots \le t_{M-1}$ of reals, all lying in some interval $[a,a+L]$, such that
\[
\nm{z_j-z_i} \le (t_j-t_i)^2 \qquad \text{for all } 0\le i,j<M.
\]
\end{definition}

If $t_i$ is literally arclength along the circle measured from a fixed point, the displayed inequality holds automatically, since chord length never exceeds arclength; so any genuine arc of length $L$ on the circle is, in particular, ordered along an arc of width $L$ in this sense, and the definition is a (mild) generalization that makes no reference to the circle's embedding in $\R^2$.

\begin{theorem}[Main theorem]\label{thm:main}
Fix $C>0$. There is $R_0=R_0(C)$ such that the following holds for every $R \ge R_0$. Let $N=R^2$ be an integer, and let $z_0,\dots,z_{M-1}$ be \emph{pairwise distinct} Gaussian integers of norm $N$, ordered along an arc of width $L \le C\sqrt{R}$. Then
\[
M \;\le\; 1168\cdot \frac{\log R}{\log\log R}.
\]
\end{theorem}

\begin{remark}
The constant $1168 = 8\times 146$ is explicit but not optimized; see Remark~\ref{rem:constants} for where it comes from. We emphasize that $R_0(C)$ is in principle computable from the proof -- nothing here is pulled from an unstated black box -- but we have not extracted an explicit value, since every step of the argument is generous with its error terms in the interest of keeping the bookkeeping tractable.
\end{remark}

\begin{remark}
The arc length scale $C\sqrt{R}$ is not an artifact of the proof technique; it is the scale at which the argument's two mechanisms (the determinant bound and the descent step) are calibrated to balance against one another, see Remark~\ref{rem:scale} after the proof of the finite form of the theorem in Section~\ref{sec:act4}. We do not know whether $O_C(\log R/\log\log R)$ continues to hold for arcs of length up to, say, $R^{1-\epsilon}$; extending the present argument to longer arcs would require re-examining the constants throughout, and may simply be false past some critical exponent, in line with the classical extremal examples for Jarn\'{i}k's theorem at length scale $\asymp R$.
\end{remark}

\begin{remark}
Compare this to the classical bound: applying Jarn\'{i}k's theorem directly to an arc of length $L=C\sqrt R$ gives only $M = O(R^{1/3})$, since $L^{2/3} \asymp R^{1/3}$. Theorem~\ref{thm:main} improves this from polynomial in $R$ to essentially logarithmic, at the cost of using the specific arithmetic of $N=R^2$ rather than the curve's convexity alone. This is consistent with -- indeed, is a strengthening at this length scale of -- the classical fact that lattice points on circles, viewed arithmetically rather than convex-geometrically, obey much stronger bounds than generic convex curves; see the discussion of Cilleruelo and C\'{o}rdoba's theorem in \cite{Chen2021}.
\end{remark}

\section{Strategy of the proof}\label{sec:strategy}

We outline the argument before carrying it out in detail, since each of the four stages below is independently routine but their interaction is the actual content of the theorem.

\textbf{Act I (Section~\ref{sec:act1}).} Fix a window size parameter $s$, to be chosen as a function of $R$ at the very end. Given $2s+1$ consecutive points $z_0,\dots,z_{2s}$ in our ordered sequence, we build a $(2s+1)\times(2s+1)$ determinant -- a disguised Vandermonde determinant in the points themselves, after using $z\bar z = N$ to clear conjugates -- whose nonvanishing is automatic (the points are distinct) and whose size, computed via the multiplicativity of the norm, is forced to be at least $N^{s^2}$. If every pairwise chord among the $2s+1$ points were shorter than some threshold $A$, the determinant would instead be bounded \emph{above} by $A^{2s(2s+1)}$ roughly; once $A$ is small enough relative to $N$, this is a contradiction, so some pair of points among any $2s+1$ consecutive ones must be farther apart than $A$. A standard sliding-window argument converts ``every block of $2s+1$ points has some far-apart pair'' into a global bound $M \ll s + sL/A$.

\textbf{Act II (Section~\ref{sec:act2}).} The threshold $A$ obtained this way is too weak by itself. We strengthen it by reducing the points modulo small primes $p$: if many points collide modulo $p$, and $p$ does not divide $N$, those collisions force extra, unexpected divisibility of the same determinant by powers of $p$ -- raising the lower bound on its size and so \emph{lowering} the threshold $A$ at which the contradiction in Act I kicks in. We quantify ``many small primes do not divide $N$'' using a fully explicit (Chebyshev-function-based) Mertens estimate, applied over a geometrically growing stack of prime intervals. Either enough primes are missing from $N$'s factorization across this stack to drive the threshold down as far as we need, or -- and this is the key dichotomy -- the Mertens estimate \emph{guarantees} that some actual prime $p$ in the stack divides $N$.

\textbf{Act III (Section~\ref{sec:act3}).} In the second case we have a genuine prime divisor $p\mid N$. Using the splitting behavior of $p$ in $\Z[i]$, every point in our configuration (if $p$ is inert) or at least half of them (if $p$ splits, by a pigeonhole argument on which of the two prime factors divides each point) can be divided by a Gaussian factor of norm $p$ or $p^2$, producing a new, smaller configuration of points on the circle of radius $\sqrt{N/p}$ or $\sqrt{N/p^2}$, with a correspondingly shorter arc. We reapply the Act~I bound to this descended configuration.

\textbf{Act IV (Section~\ref{sec:act4}).} Assembling Acts I--III gives a finite bound $M \ll s$ for an appropriately chosen window size $s$, provided $s$, the base $m_0$ and the height $k$ of the prime-interval stack are chosen to satisfy a short list of explicit inequalities. We exhibit such a choice, with $s \asymp \log R/\log\log R$, $m_0\asymp\sqrt{\log R}$ and $k\asymp \log\log R$, and verify the inequalities hold once $R$ is large enough, which finishes the proof of Theorem~\ref{thm:main}.

\section{Act I: the Ramana determinant and the basic separation bound}\label{sec:act1}

\subsection{A sliding-window lemma}

We first record the purely combinatorial fact that turns local separation into a global count; it has nothing to do with the circle and is stated for an arbitrary monotone sequence.

\begin{lemma}\label{lem:window}
Let $t_0\le t_1\le\cdots\le t_{M-1}$ all lie in an interval of length $L$. Suppose that for some integer $k\ge 1$ and real $B>0$,
\[
t_{j+k} - t_j \ge B \qquad \text{for every } 0\le j \le M-1-k.
\]
Then $M \le k + kL/B$.
\end{lemma}

\begin{proof}
Partition $\{0,\dots,M-1\}$ into the $k$ residue classes mod $k$. Within the $r$-th class, consecutive elements $t_{r}, t_{r+k}, t_{r+2k},\dots$ increase by at least $B$ at each step, so a class of size $c_r$ spans at least $(c_r-1)B \le L$, giving $c_r \le 1+L/B$. Summing over the $k$ classes, $M = \sum_r c_r \le k(1+L/B) = k+kL/B$.
\end{proof}

In every application below, $k=2s$: the relevant local statement is that any $2s+1$ \emph{consecutive} points in the sequence contain some pair at distance $\ge B$ in the $t$-parameter, which by the chord-length hypothesis is implied by some pair of squared chord-distance $\ge B^2$; we will produce such a pair from the determinant argument below.

\subsection{Ramana's determinant}

Fix $s\ge 1$ and $2s+1$ Gaussian integers $z_0,\dots,z_{2s}$, all of norm $N$. Form the $(2s+1)\times(2s+1)$ matrix
\[
\mathcal R_s(z) \;=\; \Big( \bar z_i^{\,s-j} \,\big[ j\le s\big] \;+\; z_i^{\,j-s}\,\big[j>s\big] \Big)_{0\le i,j\le 2s},
\]
that is, the $i$-th row is $(\bar z_i^{\,s},\bar z_i^{\,s-1},\dots,\bar z_i,1,z_i,z_i^2,\dots,z_i^{\,s})$, a row of both negative and positive ``powers'' of $z_i$ centered at the constant term. We call $\mathcal R_s(z)$ the \emph{Ramana matrix}; its determinant is the engine of the whole argument.

\begin{lemma}\label{lem:ramana}
With notation as above,
\[
\Big(\prod_{i=0}^{2s} z_i^{\,s}\Big)\cdot \det \mathcal R_s(z) \;=\; N^{s(s+1)/2}\cdot \prod_{0\le i<j\le 2s} (z_j-z_i).
\]
\end{lemma}

\begin{proof}[Proof sketch]
Scale row $i$ of $\mathcal R_s(z)$ by $z_i^{\,s}$. Using $\bar z_i = N/z_i$ (which holds since $z_i\bar z_i=N$), every entry $\bar z_i^{\,s-j}z_i^{\,s}$ in the left half becomes $N^{s-j}z_i^{\,j}$, and the row becomes, after factoring $N^{s-j}$ out of column $j\le s$,
\[
N^{s},\,N^{s-1}z_i,\,\dots,\,N\cdot z_i^{s-1},\,z_i^s,\,z_i^{s+1},\dots,z_i^{2s},
\]
which, dividing column $j$ by $N^{s-j}$ for $j\le s$, is exactly the $i$-th row of the ordinary $(2s+1)\times(2s+1)$ Vandermonde matrix in $z_0,\dots,z_{2s}$, raised through powers $0,1,\dots,2s$. Tracking the column scalings $N^{s},N^{s-1},\dots,N^0,1,\dots,1$ pulled out of the determinant gives the stated factor $N^{s(s+1)/2}$, and the Vandermonde determinant itself is the classical product $\prod_{i<j}(z_j-z_i)$.
\end{proof}

Taking norms on both sides of Lemma~\ref{lem:ramana}, using that $\nm{z_i}=N$ for every $i$ and that the norm is multiplicative,
\[
N^{s(2s+1)} \cdot \nm{\det\mathcal R_s(z)} \;=\; N^{s(s+1)}\cdot \prod_{i<j} \nm{z_j-z_i},
\]
and since $s(2s+1) = s(s+1)+s^2$, dividing through by $N^{s(s+1)}$ gives the basic identity we will use repeatedly:
\begin{equation}\label{eq:basic-det}
N^{s^2} \cdot \nm{\det \mathcal R_s(z)} \;=\; \prod_{0\le i<j\le 2s} \nm{z_j-z_i}.
\end{equation}

\subsection{The basic separation bound}

If $z_0,\dots,z_{2s}$ are pairwise distinct, the Vandermonde product on the right-hand side of Lemma~\ref{lem:ramana} is nonzero, hence $\det\mathcal{R}_s(z)\ne 0$ in $\Z[i]$, hence $\nm{\det \mathcal R_s(z)} \ge 1$. Plugging this trivial lower bound into \eqref{eq:basic-det} gives:

\begin{proposition}\label{prop:basic-sep}
If $z_0,\dots,z_{2s}$ are pairwise distinct Gaussian integers of norm $N$, then
\[
N^{s^2} \;\le\; \prod_{0\le i<j\le 2s} \nm{z_j-z_i}.
\]
In particular, if every pairwise squared chord length among $z_0,\dots,z_{2s}$ were at most $A^2$, the right-hand side would be at most $(A^2)^{s(2s+1)}$; so if
\begin{equation}\label{eq:subcritical}
A^{2s(2s+1)} < N^{s^2}
\end{equation}
then some pair $z_i,z_j$ among $z_0,\dots,z_{2s}$ has $\nm{z_j-z_i} > A^2$.
\end{proposition}

We call a threshold $A$ satisfying \eqref{eq:subcritical} \emph{subcritical}. Applying Proposition~\ref{prop:basic-sep} to every length-$(2s+1)$ window of consecutive points in our ordered sequence, and feeding the resulting separation into Lemma~\ref{lem:window} with $k=2s$, $B=A$:

\begin{corollary}\label{cor:basic-card}
If $A$ is subcritical for $N$ and $s$, and $z_0,\dots,z_{M-1}$ are pairwise distinct, ordered along an arc of width $L$, all of norm $N$, then
\[
M \;\le\; 2s + 2s\cdot \frac{L}{A}.
\]
\end{corollary}

Corollary~\ref{cor:basic-card} is already a real theorem -- it is, in effect, an algebraic strengthening of the classical lattice-point-on-arc bound -- but used alone it is too weak: subcriticality \eqref{eq:subcritical} forces $A \ll N^{1/(4s+2)}$, which shrinks to $1$ only once $s \gg \log N$, giving nothing better than the trivial bound $M \ll \log N$ when $L$ is comparable to $N^{O(1)}$. The point of Acts II and III is to drive $A$ down far enough, with $s$ much smaller than $\log N$, that the bound becomes nontrivial.

\section{Act II: missing primes and the Mertens dichotomy}\label{sec:act2}

\subsection{Collisions modulo a prime}

Fix a prime $p$. The reduction of the circle $x^2+y^2\equiv a \pmod p$ has at most $2p$ solutions in $(\Z/p)^2$: for each of the $p$ choices of $x$, $y^2\equiv a-x^2$ has at most $2$ solutions. Reducing our $n=2s+1$ points modulo $p$ therefore maps them into a set of size at most $2p$; if $2p$ is much smaller than $n$, the pigeonhole principle forces many \emph{collisions} -- pairs of points congruent to each other modulo $p$. Precisely, by a Cauchy--Schwarz argument on the fiber sizes of the reduction map, the number of unordered colliding pairs satisfies
\begin{equation}\label{eq:collision}
\#\{\,i<j \,:\, z_i \equiv z_j \!\!\pmod p\,\} \;\ge\; \frac{n^2}{4p} - \frac{n}{2}.
\end{equation}

\subsection{Collisions force extra divisibility}

Suppose $p \nmid N$. If $z_i \equiv z_j \pmod p$ in $\Z[i]$, the factor $p$ divides $z_j-z_i$, hence survives in the Vandermonde product $\prod_{i<j}(z_j-z_i)$ of Lemma~\ref{lem:ramana}; and since $p\nmid N$, the row-scaling factors $z_i^s$ used in the proof of that lemma are invertible modulo any power of $p$, so this extra divisibility transfers intact to $\det \mathcal R_s(z)$ itself. Writing $E_p$ for the number of colliding pairs modulo $p$, one obtains the strengthened divisibility
\[
p^{2E_p} \;\Big|\; \nm{\det \mathcal R_s(z)}.
\]
Carrying this out simultaneously over a finite set $P$ of primes not dividing $N$ -- the \emph{missing} primes -- and using that distinct primes contribute coprime factors, \eqref{eq:basic-det} upgrades to
\begin{equation}\label{eq:missing-det}
\Big(\prod_{p \in P} p^{2E_p}\Big) \cdot N^{s^2} \;\le\; \prod_{i<j} \nm{z_j-z_i}.
\end{equation}
This is exactly the same shape as Proposition~\ref{prop:basic-sep}, but with the left-hand side multiplied by an extra factor coming entirely from primes that fail to divide $N$. The more primes are missing from $N$'s factorization in a given range, the smaller a threshold $A$ can be while remaining subcritical.

To make \eqref{eq:missing-det} usable we need a clean lower bound on $E_p$ from \eqref{eq:collision}: for $p$ small relative to $s$ (specifically $4p\le s$), one has $E_p \ge \lfloor s^2/(2p)\rfloor \ge s^2/(2p) - 1$, and converting the product over $P$ into a sum of logarithms, \eqref{eq:missing-det} becomes, after a short computation absorbing the floor-loss into an additive error of size $O(\log p)$ per prime,
\begin{equation}\label{eq:log-missing}
s(2s+1)\log A^2 \;<\; \Big(s^2 \!\sum_{p\in P}\frac{\log p}{p} \;-\; 2\!\sum_{p\in P}\log p\Big) + s^2\log N
\;\Longrightarrow\; A \text{ is subcritical against } P.
\end{equation}
The sum $\sum_{p\in P} \log p / p$ appearing here is precisely the weighted prime sum $\Sigma_{\mathrm{miss}}$ of Section~\ref{sec:notation}, restricted to $P=$ the missing primes in some range; this is what ties the argument to classical Mertens-type estimates.

\subsection{Mertens estimates and the prime-interval stack}

We need a fully explicit, elementary lower bound for $\sum_{m<p\le n}\log p/p$ -- explicit enough that it can be applied at every finite scale $R$ in the final argument, not just asymptotically. Such bounds follow from Chebyshev's elementary estimates for $\theta(x)=\sum_{p\le x}\log p$ (no analytic continuation or zero-free region is needed); a short computation using $\theta(x) \asymp x$ gives, for $n\ge 8m$,
\begin{equation}\label{eq:mertens-eighth}
\sum_{m<p\le n} \frac{\log p}{p} \;\ge\; \frac{\log 2}{2},
\end{equation}
once $n$ is large enough that the (explicit) Chebyshev error terms are dominated by the main term -- a condition we verify directly in Section~\ref{sec:act4} for our particular choice of parameters. We call \eqref{eq:mertens-eighth} the \emph{eighth-scale Mertens bound}: every interval $(m,8m]$ contributes weighted prime mass at least $\log 2/2$, regardless of $m$ (once $m$ is not too small).

Stack $k$ such intervals geometrically, $I_i = (m_0\cdot 8^i,\, m_0\cdot 8^{i+1}]$ for $i=0,\dots,k-1$, all disjoint, with total range $(m_0,\,m_0\cdot 8^k]$. Summing \eqref{eq:mertens-eighth} over the stack,
\begin{equation}\label{eq:mertens-stack}
\sum_{m_0<p\le m_0\cdot 8^k} \frac{\log p}{p} \;\ge\; k\cdot \frac{\log 2}{2}.
\end{equation}

\subsection{The dichotomy}

Inequality \eqref{eq:mertens-stack} is the source of a clean either/or. Let $\Sigma_{\mathrm{miss}}(N; m_0, m_0\cdot 8^k)$ denote the weighted sum restricted to primes in the stack \emph{not} dividing $N$.

\begin{proposition}[Missing-prime dichotomy]\label{prop:dichotomy}
Exactly one of the following holds:
\begin{enumerate}[label=(\arabic*)]
\item \textbf{Many missing primes.} $\Sigma_{\mathrm{miss}}(N;m_0,m_0\cdot8^k) \ge k\log 2/2$. In this case the missing-prime mass alone may already be enough to make \eqref{eq:log-missing} hold for a threshold $A$ much smaller than the one given by Proposition~\ref{prop:basic-sep} alone.
\item \textbf{Few missing primes.} $\Sigma_{\mathrm{miss}}(N;m_0,m_0\cdot8^k) < k\log2/2$. Then, since the total weighted mass over the stack is at least $k\log2/2$ by \eqref{eq:mertens-stack} while the missing mass is strictly smaller, the \emph{present} mass -- contributed by primes that do divide $N$ -- is strictly positive: there exists a prime $p$ with $m_0 < p \le m_0\cdot 8^k$ and $p\mid N$.
\end{enumerate}
\end{proposition}

Case (1) is handled by Corollary~\ref{cor:basic-card} applied with the strengthened threshold from \eqref{eq:log-missing}, giving directly $M \le 2s + 2sL/A$. Case (2) hands us an honest prime divisor of $N$ inside the stack, and triggers Act III.

\section{Act III: prime descent}\label{sec:act3}

Suppose Proposition~\ref{prop:dichotomy} has produced a prime $p \mid N$, $p\ne 2$, with $m_0 < p$. We use the splitting of $p$ in $\Z[i]$ to shrink the problem.

\subsection{Split primes: $p \equiv 1 \pmod 4$}

By Fermat's two-square theorem, $p = \pi\bar\pi$ for some Gaussian prime $\pi$ of norm $p$. If $p \mid \nm{z}$ for a Gaussian integer $z$, then $\pi\mid z$ or $\bar\pi \mid z$ (possibly both, but not necessarily). Running over our whole configuration $z_0,\dots,z_{M-1}$, each $z_n$ is divisible by $\pi$ or by $\bar\pi$ (or both); by pigeonhole, at least half of the points are divisible by the \emph{same} one of the two factors, say $\pi$ without loss of generality. Write $z_n = \pi w_n$ for this half-subfamily. Since $\nm{z_n}=N$ and $\nm\pi = p$, we get $\nm{w_n} = N/p =: N'$: the half-subfamily $\{w_n\}$ consists of Gaussian integers of norm $N'$, i.e.\ lattice points on the smaller circle of radius $\sqrt{N'}=\sqrt{N/p}$.

Multiplication by $\pi$ scales squared distances by exactly $\nm\pi=p$: $\nm{\pi w - \pi w'} = p\cdot \nm{w-w'}$. So if the original points lay on an arc of width $L$, the descended subfamily $\{w_n\}$ lies on an arc of width at most $L/\sqrt p$ (since $t$-parameter differences scale like chord lengths, which scale like $\sqrt p$).

\subsection{Inert primes: $p \equiv 3 \pmod 4$}

If $p\equiv 3\pmod 4$, $p$ remains prime in $\Z[i]$, and the classical fact that an inert prime dividing $x^2+y^2$ must divide both $x$ and $y$ shows $p \mid z$ outright for \emph{every} point $z$ of norm divisible by $p$ -- no pigeonhole loss this time. Writing $z_n = p\, w_n$, we get $\nm{w_n} = N/p^2 =: N'$, and \emph{every} point of the original configuration descends to the circle of radius $\sqrt{N'} = \sqrt N / p$. Multiplication by $p$ scales squared distances by $p^2$, so the descended arc has width at most $L/p$.

\subsection{Reassembling the bound}

In either case, the descended subfamily is again a set of pairwise distinct Gaussian integers, ordered along a (shorter) arc, all of the same norm $N'$ -- exactly the hypotheses of Corollary~\ref{cor:basic-card}, applied now at norm $N'$ with the rescaled arc width. Working through the two cases:

\begin{proposition}[Descent bound]\label{prop:descent}
Suppose $A$ is subcritical for $N'=N/p$ (split case) or $N'=N/p^2$ (inert case) and the window parameter $s$. Then:
\[
M \le \begin{cases}
\displaystyle 2\Big(2s + 2s\cdot \frac{L/\sqrt p}{A}\Big) & p\equiv 1\!\!\pmod 4 \ \text{(split; the factor $2$ is the pigeonhole loss)},\\[2mm]
\displaystyle 2s + 2s\cdot \frac{L/p}{A} & p\equiv 3 \!\!\pmod 4 \ \text{(inert; no loss)}.
\end{cases}
\]
\end{proposition}

Since $p$ was found above the stack base, $p > m_0$, and $\sqrt{m_0}\le \sqrt p \le p$, both cases of Proposition~\ref{prop:descent} are dominated, uniformly in which prime $p$ was actually found, by the single bound
\[
M \;\le\; 2\Big(2s + 2s \cdot \frac{L/\sqrt{m_0}}{A}\Big),
\]
which depends only on the stack base $m_0$, not on the specific prime $p$. This uniform form is what we use in the final assembly.

\section{Act IV: assembling the bound}\label{sec:act4}

\subsection{The finite bound}

Combine Propositions~\ref{prop:dichotomy} and~\ref{prop:descent}. Whichever branch of the dichotomy occurs, we obtain a bound of the form $M \le 2s + 2s\cdot(\text{term})$, where $(\text{term})$ is either $L/A$ (Case 1) or $2(L/\sqrt{m_0})/A$ (Case 2, uniformly). If both terms are at most $2s$ -- a condition purely about the relative sizes of $L$, $A$, $m_0$, and $s$ -- then:

\begin{proposition}[Finite bound]\label{prop:finite}
If, with $A = C\sqrt R/\sqrt{m_0}$ (the threshold scale used at the stack's lower endpoint) and $L \le C\sqrt R$,
\begin{enumerate}[label=(\roman*)]
\item $A$ is subcritical against the missing primes in the stack $(m_0, m_0\cdot 8^k]$ at norm $N$ (so that Case 1 of Proposition~\ref{prop:dichotomy}, when it occurs, applies with this $A$), and
\item $A$ is subcritical for both descended norms $N/p$ and $N/p^2$ at the window size $s$, for every prime $m_0 < p \le m_0\cdot 8^k$,
\end{enumerate}
then
\[
M \le 8s.
\]
\end{proposition}

\begin{proof}
With $A=C\sqrt R/\sqrt{m_0}$ and $L\le C\sqrt R$, we have $L/A \le \sqrt{m_0}$ and $(L/\sqrt{m_0})/A \le 1$. Substituting these crude bounds into the two branches: Case 1 gives $M\le 2s+2s\cdot m_0$, and Case 2 gives $M \le 2(2s+2s)=8s$ outright, since $(L/\sqrt{m_0})/A \le 1$ already. Hypotheses (i) and (ii) exist precisely to license the subcriticality needed for each branch at this choice of $A$; taking the worse of the two bounds, and using $4m_0 \le s$ (part of the stack-height bookkeeping in (i)--(ii)) to control the Case-1 term, both branches collapse to $M\le 8s$.
\end{proof}

We have stated Proposition~\ref{prop:finite} somewhat loosely; the formalized version tracks the constants in (i)--(ii) exactly, in terms of an explicit relation between $s$, $m_0$, $k$, $C$ and $R$. We turn to exhibiting such a choice.

\subsection{Calibrating the parameters}

We choose
\[
s = s(R) := \big\lceil 145\cdot \tfrac{\log R}{\log\log R} \big\rceil, \qquad
m_0 = m_0(R) := \big\lfloor \sqrt{\log R}\big\rfloor, \qquad
k = k(R) := \Big\lfloor \frac{\log\log R}{8\log 8}\Big\rfloor.
\]
The window half-size $s$ is, up to the rounding, exactly of order $\log R/\log\log R$ -- which is why the final theorem has this shape, since Proposition~\ref{prop:finite} gives $M\le 8s$ outright. The stack base $m_0\asymp\sqrt{\log R}$ is large enough that the Chebyshev error terms underlying \eqref{eq:mertens-eighth} are negligible (an elementary little-$o$ estimate, since $\sqrt{\log R}\to\infty$), yet small enough that $\log m_0 \approx \frac12\log\log R$ remains a vanishing fraction of $\log N = 2\log R$ -- ensuring the subcriticality conditions (i)--(ii) of Proposition~\ref{prop:finite}, which compare $\log A$ against $\log N$ and $\log m_0$, hold with room to spare. The stack height $k\asymp\log\log R$ is calibrated so that the stack's upper endpoint $U=m_0\cdot 8^k \asymp \sqrt{\log R}\cdot(\log R)^{1/8}$ satisfies $4U\le s$ -- needed for the small-prime collision estimate $E_p\ge\lfloor s^2/2p\rfloor$ underlying \eqref{eq:log-missing} -- while still being tall enough, in units of $k$, to guarantee via \eqref{eq:mertens-stack} that the dichotomy of Proposition~\ref{prop:dichotomy} is non-vacuous.

Verifying that hypotheses (i)--(ii) of Proposition~\ref{prop:finite} hold for this choice, for all $R$ large enough (depending on $C$), is an entirely routine -- if lengthy -- exercise in elementary asymptotics ($\log\log R = o((\log R)^a)$ for every fixed $a>0$, and $(\log R)^2 = o(R)$), which we have carried out and verified mechanically; we omit the (uninstructive) details here. The upshot:

\begin{proposition}\label{prop:finite-asymp}
For every $C>0$ there is $R_0(C)$ such that for $R\ge R_0(C)$, the choice of $s,m_0,k$ above satisfies the hypotheses of Proposition~\ref{prop:finite} for every valid configuration of points with $N=R^2$ and arc width $L\le C\sqrt R$. Consequently
\[
M \le 8s(R) \le 8\cdot 146\cdot \frac{\log R}{\log\log R} = 1168\cdot\frac{\log R}{\log\log R}.
\]
\end{proposition}

(The constant $146$ rather than $145$ absorbs the ceiling in the definition of $s$.) Proposition~\ref{prop:finite-asymp} is precisely Theorem~\ref{thm:main}.

\begin{remark}\label{rem:scale}
The reason the argument is calibrated to arc width $L \asymp \sqrt R$ rather than some other power of $R$ is visible in the proof of Proposition~\ref{prop:finite}: the threshold scale $A = C\sqrt R/\sqrt{m_0}$ is forced by the descent step (a factor of $\sqrt{m_0}$ is lost passing to the smaller circle), and remaining subcritical at norm $N=R^2$ requires, very roughly, $A \ll N^{1/(4s)}$. With $s\asymp \log R/\log\log R$, this allows $A$ to be as large as $R^{O(\log\log R/\log R)} = R^{o(1)}$ -- a power of $R$ tending to $R^0$, but only just barely faster than $\sqrt{R}/\sqrt{m_0}\sim \sqrt{R}/(\log R)^{1/4}$ can afford to be. Pushing $L$ past $\sqrt R$ by a polynomial factor would force $s$ to grow polynomially as well, destroying the logarithmic shape of the final bound.
\end{remark}

\section{Remarks, constants, and outlook}\label{sec:remarks}

\begin{remark}\label{rem:constants}
We reiterate that $1168$, and the intermediate constants $145/146$, $8$, and the base-$8$ scaling of the prime-interval stack, are artifacts of a specific, formally verified but otherwise unoptimized choice of parameters. The stack base of $8$ (rather than $2$ or $e$) was chosen because it makes the eighth-scale Mertens bound \eqref{eq:mertens-eighth} comfortably true with simple, explicit error terms; a sharper Mertens estimate, or a more careful balancing of the collision-count and descent-loss terms, would very likely improve the constant $1168$ substantially, perhaps even the exponent of $\log\log R$ in the denominator. We have made no attempt at this optimization, preferring to verify a clean, correct statement first.
\end{remark}

\begin{remark}
The hypothesis in Theorem~\ref{thm:main} is stated for an abstract monotone parametrization $t$ dominating chord length, rather than for literal arclength on the circle. This is deliberate: it isolates exactly the geometric input the proof actually uses (chord $\le$ parameter-gap), and makes the theorem immediately applicable to any sequence of points satisfying this weaker hypothesis -- for instance, points ordered by argument but measured in any metric comparable to arclength. The reduction from a literal circular arc to this abstract hypothesis is immediate, since arclength dominates chord length on any circle, but we have stated the theorem in the more general form since the proof never uses anything else.
\end{remark}

\begin{remark}
We have excluded the prime $p=2$ from the descent step throughout (Section~\ref{sec:act3}), requiring every descended prime to be odd. Since $2$ ramifies, $2=-i(1+i)^2$, it could in principle be handled by the same descent mechanism using the Gaussian prime $1+i$; we have simply folded any power of $2$ dividing $N$ into the ``missing or present'' bookkeeping of the determinant branch instead, which suffices for the bound as stated and avoids a small amount of case analysis.
\end{remark}

\begin{remark}
A natural question this work does not address is what happens for arcs of length between $\sqrt R$ and $R$: does an $O_C(\log R/\log\log R)$-type bound persist, perhaps with the constant depending on the length exponent, or does the classical $R^{1/3}$-type behavior of Jarn\'{i}k's theorem reassert itself past some critical scale? We would also like to understand whether the present method, which is specific to circles via the arithmetic of $\Z[i]$, extends to lattice points on other conics -- ellipses $ax^2+by^2=N$, say, or more general binary quadratic forms -- where the relevant ring of integers is a different (possibly non-maximal, possibly non-principal) quadratic order, and the splitting behavior of primes is correspondingly more delicate. We do not address either question here.
\end{remark}

\begin{remark}
This report describes only the lattice-point bound of Theorem~\ref{thm:main}. The accompanying Lean development also contains a separate, self-contained algebraic construction -- not used anywhere in the proof above -- concerning explicit integer certificates for configurations of points in near-arithmetic-progression position on a circle, instantiated numerically on a single example built from consecutive Fibonacci numbers. Its precise relationship to the asymptotic bound of this report, if any, has not yet been worked out, and we leave its description to a future, separate note.
\end{remark}

\subsection*{Acknowledgments} This report grew out of an earlier expository project on Jarn\'{i}k's theorem and lattice points on circles \cite{Chen2021}, and the entire argument presented here has been formally verified in Lean~4 \cite{Lean4} using the Mathlib library \cite{Mathlib}.

\begin{thebibliography}{9}

\bibitem{Jarnik1926}
V.\ Jarn\'{i}k, \emph{\"{U}ber die Gitterpunkte auf konvexen Kurven}, Math.\ Z.\ \textbf{24} (1926), 500--518.

\bibitem{Chen2021}
C.\ Chen, \emph{Lattice Points On and Near Circles}, expository notes, 2021.

\bibitem{Mathlib}
The Mathlib Community, \emph{The Lean Mathematical Library}, in: Proceedings of the 9th ACM SIGPLAN International Conference on Certified Programs and Proofs, 2020. \url{https://github.com/leanprover-community/mathlib4}.

\bibitem{Lean4}
L.\ de Moura and S.\ Ullrich, \emph{The Lean 4 Theorem Prover and Programming Language}, in: Automated Deduction -- CADE 28, 2021.

\end{thebibliography}

\end{document}
